% GENERATED FILE. Edit canonical lessons, then run npm run generate:books.
% canonical-source-hash: fnv1a64:668c44764ece8d24
% canonical-lessons: GU-C422-lari, GU-C422-pheriyo, GU-C422-gallo, GU-C422-kandoi, GU-C422-khapvu

\chapter{Hawkers and Stalls}
\label{ch:gu-hawkers-and-stalls}
\hlchaptermodality{422}

\begin{chapteropening}
\textbf{By the end of this chapter:} \emph{I can buy from a street cart or a stall.}

Complete the last lesson of chapter 422: I can buy from a street cart or a stall.
\end{chapteropening}

\section[\texorpdfstring{\gu{લારી} (lārī)}{લારી (lārī)}]{\gu{લારી} (lārī) — a handcart (a street seller's)}

\label{lesson:GU-C422-lari}



\begin{quote}
Before the new one: say the Gujarati for to growl, then the Gujarati for to chase away.
\end{quote}

\subsection*{You'll want to know: \gu{લારી}}
\textbf{\gu{લારી}} — \emph{lārī} — \textquotedblleft{}a handcart (a street seller's)\textquotedblright{}.

\begin{quote}
\gu{શાકની} \gu{લારી} \gu{રોજ} \gu{અમારી} \gu{શેરીમાં} \gu{આવે} \gu{છે}.
\end{quote}

\begin{quote}
\emph{śāknī lārī roj amārī śerīmā̃ āve che.}
\end{quote}

\begin{quote}
The vegetable cart comes into our street every day.
\end{quote}

A feminine noun.

\begin{cousinweb}
\textbf{\gu{લારી}} (\emph{lārī}) No dictionary we use records an older form, so its history is not settled.
\end{cousinweb}

\begin{grammarlens}[title={What you've built}]
One more word for getting around.
\end{grammarlens}

\subsection*{Your turn}
\begin{itemize}
\raggedright
  \item \emph{Say it:} \emph{lārī}
  \item \emph{Say it:} \emph{lārī}, once more
  \item \emph{Say it:} \emph{śāknī lārī roj amārī śerīmā̃ āve che.} — The vegetable cart comes into our street every day.
  \item \emph{Recall:} say the Gujarati for to growl, then the Gujarati for to chase away, then say \emph{lārī} again
\end{itemize}

\begin{tcolorbox}[breakable,colback=teal!4,colframe=teal!35!black,title={Before you move on}]
What does \gu{લારી} mean? (\textquotedblleft{}A handcart (a street seller's)\textquotedblright{}.) Say it once more.
\end{tcolorbox}
\section[\texorpdfstring{\gu{ફેરિયો} (pheriyo)}{ફેરિયો (pheriyo)}]{\gu{ફેરિયો} (pheriyo) — a hawker; a street seller}

\label{lesson:GU-C422-pheriyo}



\begin{quote}
Before the new one: say the Gujarati for to chase away, then the Gujarati for a handcart.
\end{quote}

\subsection*{You'll want to know: \gu{ફેરિયો}}
\textbf{\gu{ફેરિયો}} — \emph{pheriyo} — \textquotedblleft{}a hawker; a street seller\textquotedblright{}.

\begin{quote}
\gu{ફેરિયો} \gu{ઘરે} \gu{ઘરે} \gu{જઈને} \gu{રમકડાં} \gu{વેચે} \gu{છે}.
\end{quote}

\begin{quote}
\emph{pheriyo ghare ghare jaīne ramakḍā̃ vece che.}
\end{quote}

\begin{quote}
The hawker goes from house to house selling toys.
\end{quote}

A masculine noun.

\begin{cousinweb}
\textbf{\gu{ફેરિયો}} (\emph{pheriyo}) is built on \gu{ફેરો}, a round, which goes back to \emph{*phēra}, \textquotedblleft{}a turn\textquotedblright{}, a form CDIAL reconstructs (CDIAL 9078): the seller who goes on rounds.
\end{cousinweb}

\begin{grammarlens}[title={What you've built}]
One more word for getting around.
\end{grammarlens}

\subsection*{Your turn}
\begin{itemize}
\raggedright
  \item \emph{Say it:} \emph{pheriyo}
  \item \emph{Say it:} \emph{pheriyo}, once more
  \item \emph{Say it:} \emph{pheriyo ghare ghare jaīne ramakḍā̃ vece che.} — The hawker goes from house to house selling toys.
  \item \emph{Recall:} say the Gujarati for to chase away, then the Gujarati for a handcart, then say \emph{pheriyo} again
\end{itemize}

\begin{tcolorbox}[breakable,colback=teal!4,colframe=teal!35!black,title={Before you move on}]
What does \gu{ફેરિયો} mean? (\textquotedblleft{}A hawker; a street seller\textquotedblright{}.) Say it once more.
\end{tcolorbox}
\section[\texorpdfstring{\gu{ગલ્લો} (gallo)}{ગલ્લો (gallo)}]{\gu{ગલ્લો} (gallo) — a paan stall; a kiosk}

\label{lesson:GU-C422-gallo}



\begin{quote}
Before the new one: say the Gujarati for a handcart, then the Gujarati for a hawker.
\end{quote}

\subsection*{You'll want to know: \gu{ગલ્લો}}
\textbf{\gu{ગલ્લો}} — \emph{gallo} — \textquotedblleft{}a paan stall; a kiosk\textquotedblright{}.

\begin{quote}
\gu{નાકા} \gu{પરના} \gu{ગલ્લા} \gu{પરથી} \gu{તેણે} \gu{પાન} \gu{લીધું}.
\end{quote}

\begin{quote}
\emph{nākā parnā gallā parthī teṇe pān līdhũ.}
\end{quote}

\begin{quote}
He bought a paan from the stall at the corner.
\end{quote}

A masculine noun.

\begin{cousinweb}
\textbf{\gu{ગલ્લો}} (\emph{gallo}) No dictionary we use records an older form, so its history is not settled.
\end{cousinweb}

\begin{grammarlens}[title={What you've built}]
One more word for getting around.
\end{grammarlens}

\subsection*{Your turn}
\begin{itemize}
\raggedright
  \item \emph{Say it:} \emph{gallo}
  \item \emph{Say it:} \emph{gallo}, once more
  \item \emph{Say it:} \emph{nākā parnā gallā parthī teṇe pān līdhũ.} — He bought a paan from the stall at the corner.
  \item \emph{Recall:} say the Gujarati for a handcart, then the Gujarati for a hawker, then say \emph{gallo} again
\end{itemize}

\begin{tcolorbox}[breakable,colback=teal!4,colframe=teal!35!black,title={Before you move on}]
What does \gu{ગલ્લો} mean? (\textquotedblleft{}A paan stall; a kiosk\textquotedblright{}.) Say it once more.
\end{tcolorbox}
\section[\texorpdfstring{\gu{કંદોઈ} (kandoī)}{કંદોઈ (kandoī)}]{\gu{કંદોઈ} (kandoī) — a sweet-maker; a sweet seller}

\label{lesson:GU-C422-kandoi}



\begin{quote}
Before the new one: say the Gujarati for a hawker, then the Gujarati for a paan stall.
\end{quote}

\subsection*{You'll want to know: \gu{કંદોઈ}}
\textbf{\gu{કંદોઈ}} — \emph{kandoī} — \textquotedblleft{}a sweet-maker; a sweet seller\textquotedblright{}.

\begin{quote}
\gu{દિવાળીમાં} \gu{કંદોઈની} \gu{દુકાને} \gu{ભીડ} \gu{છે}.
\end{quote}

\begin{quote}
\emph{divāḷīmā̃ kandoīnī dukāne bhīḍ che.}
\end{quote}

\begin{quote}
At Diwali there is a crowd at the sweet-maker's shop.
\end{quote}

A masculine noun.

\begin{cousinweb}
\textbf{\gu{કંદોઈ}} (\emph{kandoī}) is a tadbhava of Sanskrit \emph{kāndavika}, \textquotedblleft{}a baker\textquotedblright{}, built on \emph{kandu}, an iron pan; CDIAL lists the Gujarati word, \textquotedblleft{}a confectioner\textquotedblright{} (CDIAL 3034).
\end{cousinweb}

\begin{grammarlens}[title={What you've built}]
One more word for getting around.
\end{grammarlens}

\subsection*{Your turn}
\begin{itemize}
\raggedright
  \item \emph{Say it:} \emph{kandoī}
  \item \emph{Say it:} \emph{kandoī}, once more
  \item \emph{Say it:} \emph{divāḷīmā̃ kandoīnī dukāne bhīḍ che.} — At Diwali there is a crowd at the sweet-maker's shop.
  \item \emph{Recall:} say the Gujarati for a hawker, then the Gujarati for a paan stall, then say \emph{kandoī} again
\end{itemize}

\begin{tcolorbox}[breakable,colback=teal!4,colframe=teal!35!black,title={Before you move on}]
What does \gu{કંદોઈ} mean? (\textquotedblleft{}A sweet-maker; a sweet seller\textquotedblright{}.) Say it once more.
\end{tcolorbox}
\section[\texorpdfstring{\gu{ખપવું} (khapvũ)}{ખપવું (khapvũ)}]{\gu{ખપવું} (khapvũ) — to sell well; to be used up}

\label{lesson:GU-C422-khapvu}



\begin{quote}
Before the new one: say the Gujarati for a paan stall, then the Gujarati for a sweet-maker.
\end{quote}

\subsection*{You'll want to know: \gu{ખપવું}}
\textbf{\gu{ખપવું}} — \emph{khapvũ} — \textquotedblleft{}to sell well; to be used up\textquotedblright{}.

\begin{quote}
\gu{ગરમીમાં} \gu{લીંબુ} \gu{ઘણાં} \gu{ખપે} \gu{છે}.
\end{quote}

\begin{quote}
\emph{garmīmā̃ lī̃bu ghaṇā̃ khape che.}
\end{quote}

\begin{quote}
In the hot weather lemons sell well.
\end{quote}

Past (it …): \gu{તે} \gu{માલ} \gu{ખપી} \gu{ગયો}, \emph{te māl khapī gayo}.

\begin{cousinweb}
\textbf{\gu{ખપવું}} (\emph{khapvũ}) is a tadbhava of Sanskrit \emph{kṣapyatē}, \textquotedblleft{}is wasted, is destroyed\textquotedblright{}; CDIAL lists the Gujarati verb, \textquotedblleft{}to be consumed\textquotedblright{} (CDIAL 3655): goods that are used up have sold.
\end{cousinweb}

\begin{grammarlens}[title={What you've built}]
One more word for getting around.
\end{grammarlens}

\subsection*{Your turn}
\begin{itemize}
\raggedright
  \item \emph{Say it:} \emph{khapvũ}
  \item \emph{Say it:} \emph{khapvũ}, once more
  \item \emph{Say it:} \emph{garmīmā̃ lī̃bu ghaṇā̃ khape che.} — In the hot weather lemons sell well.
  \item \emph{Recall:} say the Gujarati for a paan stall, then the Gujarati for a sweet-maker, then say \emph{khapvũ} again
\end{itemize}

\begin{tcolorbox}[breakable,colback=teal!4,colframe=teal!35!black,title={Before you move on}]
What does \gu{ખપવું} mean? (\textquotedblleft{}To sell well; to be used up\textquotedblright{}.) Say it once more.
\end{tcolorbox}
